%-------------------------------------------------------------------------------
%	ABSTRACT PAGE
%-------------------------------------------------------------------------------

\addchaptertocentry{Хураангуй} % Хураангуйг гарчигт нэмэх

\begin{center}
{\scshape\Large \univname\par} % Их сургуулийн нэр
{\scshape\large \facname\par}\vspace{0.5cm} % Их сургуулийн нэр
{\huge\textbf{{Хураангуй}} \par}
\bigskip
{\Large{\ttitle} \par} % Тезисийн нэр
\bigskip

{\normalsize \shortname \par} % Зохиогчийн нэр
\addressname
\end{center}
\bigskip
\qquad Агуулахын бараа материал, түүхий эд, хангамжийн материал, үндсэн
хөрөнгө зэрэг хөрөнгүүдийг татан авах, хадгалах, дамжуулах, зарлагадах зэрэг үйл ажиллагаанд зориулж агуулахын удирдлагын системийг хөгжүүлсэн. Агуулахын үйл ажиллагаатай холбоотой өгөгдөл болон програмуудыг судалж энэхүү бичиг баримтыг боловсруулсан. Байгууллагын үйлдвэрлэл, борлуулалт, төлөвлөлтийн үйл ажиллагааг тасралтгүй байлгаж, улмаар үргүй зардал гаргахгүйн тулд Агуулахын төлөвлөлт зохион байгуулалтыг хэрхэн зөв зохистой хурдан явуулах, хэрхэн барааны  дуусах хугацаа нь дөхөж байгаа барааны мэдээллийг автоматаар болон жагсаалтаар санал болгох, барааны оруулах гаргах хүсэлтээр агуулахын барааны хөдөлгөөнийг хэрхэн удирдах системийн өгөгдлийн санг ашиглан тайлан мэдээллийг олон загвараар шууд гаргах боломжтой болгох зэрэг аргыг судалсан.



